\documentclass[letterpaper]{article}
\usepackage[submission]{aaai2027}
\usepackage[hyphens]{url}
\usepackage{graphicx}
\urlstyle{rm}
\def\UrlFont{\rm}
\usepackage{natbib}
\usepackage{caption}
\usepackage{booktabs}
\frenchspacing
\pdfinfo{/TemplateVersion (2027.1)}
\setcounter{secnumdepth}{0}
\title{Counterfactual Calendars for Diagnosing Multi-Agent Scheduling}
\author{Anonymous Submission}
\affiliations{}
\begin{document}
\maketitle

\begin{abstract}

We audit 200 MASBench scheduling cases and construct 148 counterfactual calendar pairs with disjoint feasible meeting sets. The pairs allow distinguishing constraints to be moved between participants while preserving each world's answer set. Using a subset of these pairs, we test whether changing the order of two candidate lines affects a scheduling leader's decision. Each decision contains explicit availability statements for five participants and exactly one feasible candidate. On 32 held-out scenarios, each evaluated in two availability worlds, Qwen3-8B with 4-bit BitsAndBytes quantization selected the feasible candidate in 56 of 64 worlds when candidate A was listed first and 64 of 64 when candidate B was listed first. All eight changed answers went from an incorrect A to a correct B. The paired accuracy difference was 12.5 percentage points, with a scenario-clustered bootstrap interval of 4.7 to 20.3 points.

\end{abstract}
\section{Introduction}

A scheduler that aggregates participant constraints should choose a meeting everyone can attend. A leader may receive the same participant information in different layouts, and the display of candidates may influence its decision even when availability is explicit. We hold participant reports fixed to test the leader's decision separately from communication failures, using deterministic intersection as a reference.

We investigate a controlled decision derived from MASBench scheduling cases. Each prompt presents two candidate times and five participant messages. Exactly one candidate is available to everyone. We compare two prompts for the same decision that differ only in the order of the candidate lines. The candidate labels remain attached to the same times; the participant messages continue to describe A before B.

Our calendar audit establishes which cases support moving distinguishing constraints between participants while preserving the feasible meeting sets. We use the resulting pairs to isolate the leader's decision with explicit availability messages. On a separate set of 32 scenarios, reversing the candidate lines corrected eight decisions and worsened none under the tested configuration.

\section{Related work}

Multi-agent debate studies how agents exchange candidate answers and reasoning over multiple rounds \cite{du2024debate}. MASBench studies multi-agent collaboration under partial observability, including scheduling with private participant calendars \cite{chu2026masbench}. We use its scheduling calendars to construct explicit leader inputs. The participant messages are generated from calendar availability, not obtained from a model conversation. Prior work reports substantial option-order sensitivity in multiple-choice evaluation \cite{pezeshkpour2024order}, while other work studies sensitivity to option symbols \cite{yang2025symbol}. Zheng et al. trace selection bias in multiple-choice tasks to prior preferences for option-ID tokens \cite{zheng2024selectors}.

\section{Task and experimental design}

\subsection{Calendar audit}

We audited 200 scheduling cases. In 189 cases, inserted busy entries could be identified in one participant's calendar. Removing these entries expanded the feasible meeting set. Replaying the entries into each of the other four calendars preserved the released answer set. For 148 cases, we also constructed alternative blocks yielding an equally sized, disjoint feasible set. The audit checked that moving the blocks to another participant preserved the answer sets in both worlds. These pairs enable controlled studies of where distinguishing evidence is held; the selector experiment below uses the original placement.

\subsection{Calendar pairs and candidate selection}

Each scenario contains released and alternative calendars. The two worlds have disjoint feasible meeting starts. One designated participant's calendar changes between worlds, while the other four participants' availability remains fixed. From each pair, we deterministically select one 20-minute meeting feasible in the released world and one feasible in the alternative world. We require the designated participant's availability for those candidates to reverse between worlds and all four other participants to be available for both. The saved calendar-derived inputs were checked against the paired plan. Candidate A is the released-world candidate for even-indexed eligible pairs and the alternative-world candidate for odd-indexed pairs; B receives the other candidate.

Only one participant discriminates between candidates; the other four are available for both. The task therefore does not require combining multiple independently necessary constraints. Each world is rendered as two labeled candidate lines followed by five factual messages of the form \texttt{Member1: A available; B unavailable.} A deterministic intersection of the availability statements yields exactly one correct candidate. The leader is asked to output exactly \texttt{FINAL A} or \texttt{FINAL B}. Each world is evaluated independently; the model is not told the answer in the other world.

\subsection{Example}
One A-first prompt read:
\begin{quote}\small
Choose the one meeting that all five participants can attend.\\
A: Day 2 15:50-16:10\\
B: Day 1 20:40-21:00\\
Participant messages:\\
Member1: A unavailable; B available.\\
Member2: A available; B available.\\
Member3: A available; B available.\\
Member4: A available; B available.\\
Member5: A available; B available.\\
Output exactly one line: FINAL A or FINAL B.
\end{quote}
The B-first prompt exchanges the candidate lines. The model returned \texttt{FINAL A} for A-first and \texttt{FINAL B} for B-first; B is feasible.

\subsection{Listing-order intervention}

The A-first prompt lists the line for A followed by the line for B. The B-first prompt exchanges those two lines. Their labels, candidate times, availability statements, message order, and output instruction remain unchanged. The intervention changes the position of both candidate lines while leaving other A-before-B references intact. It does not distinguish all possible effects of labels, position, and distance from subsequent text.

\subsection{Data and evaluation}

Of the 148 audited pairs, 123 had no unexplained busy units and qualified for reconstruction. All 87 cases outside the pilot and exploratory sets met the candidate rule. We selected 32 by stratified sampling using seed 9401. The confirmation set has 64 world decisions, each evaluated under both listing orders, for 128 model calls. We chose the comparison after inspecting exploratory results; the confirmation scenarios were held out from that development.

The primary outcome is exact-format decision accuracy per world. Any output other than the required single line is scored as incorrect. We report paired outcomes and the B-first minus A-first accuracy difference. A percentile bootstrap resamples scenarios, retaining both worlds and both prompt orders together, for 10,000 draws with seed 6975.

The main model was Qwen3-8B served by vLLM with BitsAndBytes 4-bit quantization, temperature zero, thinking disabled, and a 256-token output limit. Calls for the two orders were interleaved by scenario and world. The server was configured with revision \texttt{main}.

\section{Results}

\subsection{Held-out listing-order comparison}

\begin{table}[t]
\centering
\begin{tabular}{lrr}
\toprule
A-first & B-first correct & B-first incorrect \\
\midrule
Correct & 56 & 0 \\
Incorrect & 8 & 0 \\
\bottomrule
\end{tabular}
\caption{Paired outcomes across 64 held-out world decisions.}
\end{table}

A-first accuracy was 56/64 (87.5\%); B-first accuracy was 64/64 (100\%). The paired difference was 12.5 percentage points, with a scenario-bootstrap 95\% interval of 4.7 to 20.3 percentage points. All eight changed answers moved from \texttt{FINAL A} to the correct \texttt{FINAL B}. The remaining 56 outputs were unchanged and correct. For A-correct worlds accuracy was 32/32 under both orders; for B-correct worlds it was 24/32 with A first and 32/32 with B first. At the scenario level, both worlds were correct in 24/32 scenarios with A first and 32/32 with B first (two-sided exact sign test on discordant scenarios, $p=0.0078$). By evidence holder, the numbers of corrected decisions out of all decisions were 4/20 for Member1, 1/6 for Member2, 2/14 for Member3, 1/10 for Member4, and 0/14 for Member5. All outputs met the strict format.

\subsection{Exploratory observations}

In the preceding 32-scenario exploratory set, the same listing-line reversal changed ten incorrect A responses to correct B responses, with no regressions. In a separate candidate-label swap, the model chose A in both presentations in 21 of 64 paired world comparisons, thereby changing the chosen meeting when the labels were exchanged. No comparison retained B in both presentations. The swap couples label and candidate position, so these counts cannot isolate a token preference. The exploratory analyses motivated the held-out comparison.

\section{Discussion}

Changing the order of the candidate listing lines changed Qwen3-8B's decisions on eight held-out worlds. The intervention shifts line position and proximity to the participant messages together.

The calendar pairs support comparisons across evidence placements with the same feasible meetings. The demonstration tests the leader with fixed, correct reports and supplied candidates. Candidate discovery, communication, and negotiation remain outside its scope.

The sample consists of 32 selected scenarios from one benchmark and one prompt family. The main result uses one model and quantization configuration at temperature zero. A 4B run failed the strict output format; exploratory label extraction is reported in the supplement. The startup log records a branch name, and the exact checkpoint commit is unknown.

\section{Conclusion}

The calendar audit provides 148 paired scheduling cases for controlled studies of evidence allocation. In the leader demonstration, reversing two candidate listing lines changed eight of 64 held-out decisions for the tested Qwen3-8B configuration, despite identical availability evidence. Strict accuracy improved from 87.5\% to 100\% on this sample.

\section{AI use}
Generative AI assisted with manuscript editing and analysis code.

\bibliography{references}
\end{document}
